\Needspace{6\baselineskip}
\begin{center}\hypertarget{native-257-article184}{}\pdfbookmark[2]{184.}{article-184}184.\end{center}
Um die vorangegangenen Untersuchungen an einem Beispiel zu veranschaulichen, nehmen wir an, dass Beobachtungen von vermutlich gleicher Präzision zu den Gleichungen
\[
\begin{aligned}
p-q+2r&=3\\
3p+2q-5r&=5\\
4p+q+4r&=21,
\end{aligned}
\]
geführt haben und dass \SourceLeaf{258}eine vierte Beobachtung, der nur die halbe Präzision zuzuschreiben ist, die Gleichung
\[
-2p+6q+6r=28.
\]
ergeben hat. Wir ersetzen daher die letzte Gleichung durch
\[
-p+3q+3r=14
\]
und nehmen an, dass sie aus einer Beobachtung von derselben Präzision wie die vorangegangenen hervorgegangen ist. Damit ergibt sich
\[
\begin{aligned}
P&=27p+6q-88\\
Q&=6p+15q+r-70\\
R&=q+54r-107,
\end{aligned}
\]
und daraus durch Elimination
\[
\begin{aligned}
19899p&=49154+809P-324Q+6R\\
737q&=2617-12P+54Q-R\\
6633r&=12707+2P-9Q+123R.
\end{aligned}
\]
Die wahrscheinlichsten Werte der Unbekannten sind also
\[
\begin{aligned}
p&=2{,}470\\
q&=3{,}551\\
r&=1{,}916,
\end{aligned}
\]
und die diesen Bestimmungen zuzuschreibende relative Präzision beträgt, wenn man die Präzision der ursprünglichen Beobachtungen $=1$ setzt,
\[
\begin{aligned}
\text{für }p\quad .\,.\,.\,.\quad &\BareRoot{\frac{19899}{809}}=4{,}96\\[4pt]
\text{für }q\quad .\,.\,.\,.\quad &\BareRoot{\frac{737}{54}}=3{,}69\\[4pt]
\text{für }r\quad .\,.\,.\,.\quad &\BareRoot{\frac{2211}{41}}=7{,}34.
\end{aligned}
\]
